\pageoriginale

\clearpage
\SGAExposeBookmark{XXII}{Una fórmula de congruencia para la función zeta}
\begin{center}
{\large SGA 7 --- EXPOSÉ XXII}\par
\vspace{2em}
{\Large\underline{Una fórmula de congruencia para la función zeta}}\par
\vspace{1.5em}
{\large por N. Katz}\par
\end{center}

\section*{0. Introducción}

Sea $X$ un esquema propio sobre $\mathbb F_q$. Su función zeta,
considerada como elemento de $\mathbb Z[[T]]$, puede reducirse módulo $p$;
queremos calcularla módulo $p$. Obtenemos una expresión del tipo que cabía
esperar de las fórmulas de puntos fijos de tipo «Woods Hole»: un producto
alternado de polinomios característicos de un «Frobenius» adecuado que
actúa sobre cualquiera de los tres espacios de cohomología siguientes:

\begin{enumerate}[label=\arabic*)]
\item la cohomología coherente $H^i(X,\mathcal O_X)$ (3.1.1);
\item la cohomología de De Rham
  $H^i_{\mathrm{DR}}(X)=\mathbb H^i(X,\Omega_X^\bullet)$ (3.1.1);
\item la cohomología étale $H^i_{\mathrm{et}}(X,\mathbb Z/p\mathbb Z)$
  (3.1.2).
\end{enumerate}

Esta expresión debe verse como un perfeccionamiento de un caso particular
de la fórmula de trazas de Woods Hole, a saber, el del haz estructural y el
endomorfismo de Frobenius. En efecto, este caso es un corolario de la
fórmula de congruencia (cf.~3.2.1).

El método de esta exposición es completamente ingenuo. Consiste en reducirse al
caso de una hipersuperficie y después «reducir módulo $p$» el cálculo de
Dwork para la verdadera función zeta de una hipersuperficie.

\pageoriginale

Al efectuar cuidadosamente esta «reducción», se obtiene el
perfeccionamiento 5.0.6, debido a Ax~[1] (y también a Dwork en el caso liso
[4, \S~7]), del conocido teorema de Warning--Chevalley~[8].

Deligne acaba de encontrar otro método, enteramente cohomológico y en el
espíritu de SGA~5, que se aplica también a los «coeficientes torcidos» y que
debería aparecer en la nueva edición de SGA~5. Se apoya en la fórmula de
trazas de tipo Woods Hole y en la «fórmula de Künneth simétrica» de Deligne
(SGA~4, XVII~5.5.21).

\section*{\texorpdfstring{1. Recordatorio sobre las operaciones
$p$-lineales}{1. Recordatorio sobre las operaciones p-lineales}}

\subsection*{1.0.}
Sea $V$ un espacio vectorial de dimensión finita sobre un cuerpo $k$ de
característica $p>0$, y sea $q$ una potencia de $p$. Una aplicación

\begin{equation}\tag{1.0.1}
\varphi:V\longrightarrow V
\end{equation}

se llama $q$-\emph{lineal} si es aditiva y satisface

\begin{equation}\tag{1.0.2}
\varphi(\lambda v)=\lambda^q\varphi(v)
\qquad\text{para todo }\lambda\in k,\ v\in V.
\end{equation}

\subsection*{1.0.3.}
Obsérvese que la $n$-ésima iterada $\varphi^n$ de un endomorfismo
$q$-lineal es $q^n$-lineal y que, si $k$ no es perfecto, la imagen
$\varphi(V)$ de un endomorfismo $q$-lineal no es necesariamente un
subespacio vectorial sobre $k$.

\subsection*{1.0.4.}
Para toda extensión $K/k$, un endomorfismo $q$-lineal $\varphi$ de $V$
da lugar a un único endomorfismo $q$-lineal $\varphi_K$ de
$V_K=V\otimes_kK$ que satisface

\begin{equation}\tag{1.0.5}
\varphi_K(\mu\otimes v)=\mu^q\varphi(v)
\qquad\text{para }\mu\in K,\ v\in V.
\end{equation}

\pageoriginale

\subsection*{1.0.6.}
Un endomorfismo $q$-lineal $\varphi$ de $V$ se llama \emph{semisimple} si
su imagen genera $V$ sobre $k$, y \emph{nilpotente} si alguna de sus
iteradas es cero. Definamos los siguientes subespacios de $V$, estables por
$\varphi$:

\begin{equation}\tag{1.0.7}
V_{\mathrm{ss}}=
\bigcap_{n\geq1}\langle\varphi^n(V)\rangle_k,
\end{equation}

\begin{equation}\tag{1.0.8}
V_{\mathrm{nilp}}=
\bigcup_{n\geq1}\ker(\varphi^n:V\longrightarrow V).
\end{equation}

Es claro que la restricción de $\varphi$ a $V_{\mathrm{ss}}$ es
semisimple, mientras que su restricción a $V_{\mathrm{nilp}}$ es
nilpotente. Si $k$ es perfecto, entonces

\begin{equation}\tag{1.0.9}
V=V_{\mathrm{ss}}\oplus V_{\mathrm{nilp}}.
\end{equation}

Si $k$ no es perfecto, en cambio, (1.0.9) es falsa, pues existen
endomorfismos $q$-lineales inyectivos que no son semisimples. Para una
extensión arbitraria $K/k$, se tiene

\begin{equation}\tag{1.0.10}
(V_K)_{\mathrm{ss}}=V_{\mathrm{ss}}\otimes_kK.
\end{equation}

\subsubsection*{Proposición 1.1.}
Sea $\varphi$ un endomorfismo $q$-lineal de un espacio vectorial $V$ de
dimensión finita sobre un cuerpo separablemente cerrado $k$. Supongamos que
$\varphi$ es semisimple, y pongamos

\[
V^{1-\varphi}
=\text{el espacio vectorial sobre $\mathbb F_q$ de los puntos fijos de
$\varphi$}
=\ker(1-\varphi).
\]

Entonces

\begin{equation}\tag{1.1.1}
V\xleftarrow{\ \sim\ }
V^{1-\varphi}\otimes_{\mathbb F_q}k.
\end{equation}

\pageoriginale

\subsubsection*{Corolario 1.1.2.}
Hay un isomorfismo

\[
V^{1-\varphi^n}\xleftarrow{\ \sim\ }
V^{1-\varphi}\otimes_{\mathbb F_q}\mathbb F_{q^n}.
\]

\medskip
\noindent\emph{Demostración de la Proposición 1.1.}
(Lang, cf.~[6], p.~119.) Sea $e$ una base de $V$ sobre $k$, y escribamos

\begin{equation}\tag{1.1.3}
\varphi(e)=Ae,\qquad A\in\mathrm{GL}(n,k).
\end{equation}

Sea

\begin{equation}\tag{1.1.4}
F:\mathrm{GL}(n,k)\longrightarrow\mathrm{GL}(n,k)
\end{equation}

la aplicación que eleva cada coordenada a la $q$-ésima potencia. Para
$g\in\mathrm{GL}(n,k)$, la base $ge$ de $V$ está formada por puntos fijos
de $\varphi$ si y solo si

\begin{equation}\tag{1.1.5}
F(g)Ag^{-1}=1.
\end{equation}

Sea $\bar k$ una clausura algebraica de $k$. El grupo
$\mathrm{GL}(n,\bar k)$ actúa sobre la variedad
$\mathrm{GL}(n,\bar k)$ mediante

\begin{equation}\tag{1.1.6}
\left\{
\begin{aligned}
\mathrm{GL}(n,\bar k)\times\mathrm{GL}(n,\bar k)
&\longrightarrow\mathrm{GL}(n,\bar k),\\
(g,A)&\longmapsto F(g)Ag^{-1}.
\end{aligned}
\right.
\end{equation}

Para $A$ fijo, la aplicación

\begin{equation}\tag{1.1.7}
g\longmapsto F(g)Ag^{-1}
\end{equation}

es étale, como se ve calculando su aplicación tangente. Sus órbitas son,
por tanto, abiertas. Como $\mathrm{GL}(n,\bar k)$ es irreducible, no hay
más que una órbita; en otras palabras, la aplicación

\begin{equation}\tag{1.1.8}
\begin{aligned}
\mathrm{GL}(n,\bar k)&\longrightarrow\mathrm{GL}(n,\bar k),\\
g&\longmapsto F(g^{-1})\,g
\end{aligned}
\end{equation}

\pageoriginale

es sobreyectiva. Puesto que es además étale y $k$ es separablemente
cerrado, la aplicación

\begin{equation}\tag{1.1.9}
\begin{aligned}
\mathrm{GL}(n,k)&\longrightarrow\mathrm{GL}(n,k),\\
g&\longmapsto F(g^{-1})\,g
\end{aligned}
\end{equation}

es sobreyectiva. Así, $V$ posee una base sobre $k$ formada por puntos fijos
de $\varphi$. Como $\varphi$ es $q$-lineal, todo punto fijo de $\varphi$
pertenece al espacio vectorial sobre $\mathbb F_q$ generado por una tal base
sobre $k$. Esto demuestra la Proposición~1.1.

\subsubsection*{Corolario 1.1.10.}
Sea $\varphi$ un endomorfismo $q$-lineal de un espacio vectorial $V$ de
dimensión finita sobre un cuerpo separablemente cerrado $k$. Hay un
isomorfismo canónico

\begin{equation}\tag{1.1.11}
V_{\mathrm{ss}}\xleftarrow{\ \sim\ }
V^{1-\varphi}\otimes_{\mathbb F_q}k.
\end{equation}

En efecto, basta aplicar la Proposición~1.1 a $V_{\mathrm{ss}}$.

\subsubsection*{Proposición 1.2.}
Sea $\varphi$ un endomorfismo $q$-lineal de un espacio vectorial de
dimensión finita sobre un cuerpo separablemente cerrado $k$. Entonces la
aplicación aditiva

\[
1-\varphi:V\longrightarrow V
\]

es sobreyectiva.

\medskip
\noindent\emph{Demostración.}
Hay una sucesión exacta de grupos abelianos

\[
0\longrightarrow V_{\mathrm{ss}}\longrightarrow V
\longrightarrow V/V_{\mathrm{ss}}\longrightarrow0,
\]

\pageoriginale

sobre la cual actúa $1-\varphi$. Por el lema de la serpiente, basta tratar
$V_{\mathrm{ss}}$, lo cual es inmediato por la Proposición~1.1 usando una
base formada por puntos fijos, y $V/V_{\mathrm{ss}}$, sobre el cual
$\varphi$ es nilpotente y, por tanto, $1-\varphi$ es invertible.

\section*{2. Cohomología coherente y cohomología étale}

\subsubsection*{Proposición 2.0.}
Sea $\bar X$ un esquema propio sobre un cuerpo separablemente cerrado $k$,
y sea

\begin{equation}\tag{2.0.1}
\mathfrak F_p:H^i(\bar X,\mathcal O_{\bar X})
\longrightarrow H^i(\bar X,\mathcal O_{\bar X})
\end{equation}

la aplicación $p$-lineal inducida por

\begin{equation}\tag{2.0.2}
\mathfrak F_p:\mathcal O_{\bar X}\longrightarrow\mathcal O_{\bar X},
\qquad \mathfrak F_p(f)=f^p.
\end{equation}

Denotemos por

\[
H^i_{\mathrm{et}}(\bar X,\mathbb Z/p\mathbb Z)
\]

los grupos de cohomología étale con coeficientes en el haz constante
$\mathbb Z/p\mathbb Z$. Entonces hay isomorfismos canónicos

\begin{equation}\tag{2.0.3}
H^i_{\mathrm{et}}(\bar X,\mathbb Z/p\mathbb Z)
\simeq
H^i(\bar X,\mathcal O_{\bar X})^{1-\mathfrak F_p}.
\end{equation}

\medskip
\noindent\emph{Demostración.}
Hay una sucesión exacta de haces sobre $\bar X_{\mathrm{et}}$:

\begin{equation}\tag{2.0.4}
0\longrightarrow\mathbb Z/p\mathbb Z\longrightarrow\mathcal O_{\bar X}
\xrightarrow{\ 1-\mathfrak F_p\ }\mathcal O_{\bar X}\longrightarrow0,
\end{equation}

que da una sucesión exacta de cohomología

\begin{equation}\tag{2.0.5}
\cdots\longrightarrow H^i_{\mathrm{et}}(\bar X,\mathbb Z/p\mathbb Z)
\longrightarrow H^i_{\mathrm{et}}(\bar X,\mathcal O_{\bar X})
\xrightarrow{\ 1-\mathfrak F_p\ }
H^i_{\mathrm{et}}(\bar X,\mathcal O_{\bar X})\longrightarrow\cdots .
\end{equation}

Se tiene
$H^i_{\mathrm{et}}(\bar X,\mathcal O_{\bar X})\simeq
H^i(\bar X,\mathcal O_{\bar X})$; este es un espacio vectorial de
dimensión finita sobre $k$ porque $\bar X$ es propio. Aplicando la
Proposición~1.2 se obtienen las sucesiones exactas

\pageoriginale

\begin{equation}\tag{2.0.6}
0\longrightarrow H^i_{\mathrm{et}}(\bar X,\mathbb Z/p\mathbb Z)
\longrightarrow H^i(\bar X,\mathcal O_{\bar X})
\xrightarrow{\ 1-\mathfrak F_p\ }
H^i(\bar X,\mathcal O_{\bar X})\longrightarrow0,
\end{equation}

lo que demuestra la afirmación. La Proposición~1.1.10 da ahora el
corolario siguiente de la Proposición~2.0.

\subsubsection*{Corolario 2.1.}
Para $\bar X$ propio sobre un cuerpo separablemente cerrado $k$, hay un
isomorfismo canónico

\begin{equation}\tag{2.1.1}
H^i(\bar X,\mathcal O_{\bar X})_{\mathrm{ss}}
\xleftarrow{\ \sim\ }
H^i_{\mathrm{et}}(\bar X,\mathbb Z/p\mathbb Z)
\otimes_{\mathbb F_p}k.
\end{equation}

\subsection*{2.2.}
Sea ahora $f:X\longrightarrow\operatorname{Spec}(\mathbb F_q)$ un esquema propio sobre
$\mathbb F_q$. Sea $k$ una clausura algebraica de $\mathbb F_q$, y
pongamos

\begin{equation}\tag{2.2.1}
\bar X=X\times_{\mathbb F_q}k.
\end{equation}

El haz étale $R^if_*(\mathbb Z/p\mathbb Z)$ sobre
$\operatorname{Spec}(\mathbb F_q)$ se identifica con
$H^i_{\mathrm{et}}(\bar X,\mathbb Z/p\mathbb Z)$, considerado como
$\operatorname{Gal}(k/\mathbb F_q)$-módulo. Sea

\begin{equation}\tag{2.2.2}
\mathfrak F_q:H^i(X,\mathcal O_X)\longrightarrow H^i(X,\mathcal O_X)
\end{equation}

la operación lineal (es $q$-lineal, pero el cuerpo base es
$\mathbb F_q$) inducida por

\begin{equation}\tag{2.2.3}
\mathfrak F_q:\mathcal O_X\longrightarrow\mathcal O_X,
\qquad \mathfrak F_q(f)=f^q.
\end{equation}

\subsubsection*{Proposición 2.2.4.}
Sea $\varphi_q\in\operatorname{Gal}(k/\mathbb F_q)$ el elemento «Frobenius»
$\varphi_q(\lambda)=\lambda^q$, y sea $i\in\mathbb Z$. Hay una identidad
de polinomios característicos

\begin{equation}\tag{2.2.4.1}
\det_{\mathbb F_p}\!\left(
1-t\varphi_q^{-1}\mid H^i_{\mathrm{et}}(\bar X,\mathbb Z/p\mathbb Z)
\right)
=
\det_{\mathbb F_q}\!\left(
1-t\mathfrak F_q\mid H^i(X,\mathcal O_X)
\right).
\end{equation}

Aquí, el miembro de la izquierda es el polinomio característico del
elemento $\varphi_q^{-1}$ de $\pi_1(\operatorname{Spec}(\mathbb F_q))$ que actúa sobre
el módulo de Galois
$H^i_{\mathrm{et}}(\bar X,\mathbb Z/p\mathbb Z)$.

La Proposición~2.2.4 es un corolario del resultado siguiente.

\pageoriginale

\subsubsection*{Proposición 2.2.5.}
Sea $X$ propio sobre $\operatorname{Spec}(\mathbb F_q)$. Hay un isomorfismo canónico

\begin{equation}\tag{2.2.6}
H^i(X,\mathcal O_X)_{\mathrm{ss}}\otimes_{\mathbb F_q}k
\simeq
H^i_{\mathrm{et}}(\bar X,\mathbb Z/p\mathbb Z)
\otimes_{\mathbb F_p}k
\end{equation}

que lleva el automorfismo inducido por $\mathfrak F_q$ en el miembro de la
izquierda al automorfismo inducido por $\varphi_q^{-1}$ en el miembro de la
derecha. Este último puede interpretarse también como
$(\mathrm{Fr}_q^X)^*\otimes_{\mathbb F_p}k$, donde
$\mathrm{Fr}_q^X:\bar X\longrightarrow\bar X$ es el endomorfismo de
Frobenius sobre $k$ (cf.~SGA~5, XVIII~2.3.4).

\medskip
\noindent\emph{Demostración.}
Puesto que

\begin{equation}\tag{2.2.7}
H^i(\bar X,\mathcal O_{\bar X})\simeq
H^i(X,\mathcal O_X)\otimes_{\mathbb F_q}k,
\end{equation}

la Proposición~1.0.10 da

\begin{equation}\tag{2.2.8}
H^i(\bar X,\mathcal O_{\bar X})_{\mathrm{ss}}\simeq
H^i(X,\mathcal O_X)_{\mathrm{ss}}\otimes_{\mathbb F_q}k.
\end{equation}

Así, el isomorfismo (2.2.6) se obtiene componiendo (2.2.8) con (2.1.1).

En virtud de (2.2.7), la operación $q$-lineal

\begin{equation}\tag{2.2.9}
\mathfrak F_q:H^i(X,\mathcal O_X)\longrightarrow H^i(X,\mathcal O_X)
\end{equation}

posee una extensión $q$-lineal (cf.~1.0.5)

\begin{equation}\tag{2.2.10}
\overline{\mathfrak F}_q:H^i(\bar X,\mathcal O_{\bar X})
\longrightarrow H^i(\bar X,\mathcal O_{\bar X}),
\end{equation}

y también una extensión $k$-lineal, denotada de nuevo por
$\overline{\mathfrak F}_q$,

\begin{equation}\tag{2.2.11}
\overline{\mathfrak F}_q:H^i(\bar X,\mathcal O_{\bar X})
\longrightarrow H^i(\bar X,\mathcal O_{\bar X}).
\end{equation}

Sea

\pageoriginale

\begin{equation}\tag{2.2.12}
\varphi_q:H^i(\bar X,\mathcal O_{\bar X})
\longrightarrow H^i(\bar X,\mathcal O_{\bar X})
\end{equation}

la operación $q$-lineal cuyos puntos fijos son
$H^i(X,\mathcal O_X)$. Bajo (2.2.7), los operadores $\mathfrak F_q$,
$\varphi_q$ y $\overline{\mathfrak F}_q$ pueden escribirse también,
respectivamente, como $\mathfrak F_q\otimes\mathrm{id}_k$,
$\mathrm{id}\otimes\varphi_q$ y
$\mathfrak F_q\otimes\varphi_q$. Por tanto, conmutan entre sí y

\begin{equation}\tag{2.2.13}
\overline{\mathfrak F}_q
=\varphi_q\mathfrak F_q
=\mathfrak F_q\circ\varphi_q.
\end{equation}

Todos conmutan con $\mathfrak F_p$, por lo que $\mathfrak F_q$ y
$\varphi_q$ actúan tanto sobre
$H^i(\bar X,\mathcal O_{\bar X})^{1-\overline{\mathfrak F}_q}$ como sobre
$H^i(\bar X,\mathcal O_{\bar X})^{1-\mathfrak F_p}$. Por (2.2.13),
$\mathfrak F_q$ y $\varphi_q$ son inversos el uno del otro sobre
$H^i(\bar X,\mathcal O_{\bar X})^{1-\overline{\mathfrak F}_q}$. De ahí
que la acción de $\mathfrak F_q$ sobre

\[
H^i(\bar X,\mathcal O_{\bar X})_{\mathrm{ss}}
\simeq
H^i(\bar X,\mathcal O_{\bar X})^{1-\overline{\mathfrak F}_q}
\otimes_{\mathbb F_q}k
\qquad\text{(cf.~1.1.11)}
\]

sea inversa a la acción de $\varphi_q\otimes\mathrm{id}_k$ sobre

\[
H^i(\bar X,\mathcal O_{\bar X})^{1-\overline{\mathfrak F}_q}
\otimes_{\mathbb F_q}k.
\]

Escribiendo

\begin{equation}\tag{2.2.14}
H^i(\bar X,\mathcal O_{\bar X})^{1-\overline{\mathfrak F}_q}
\simeq
H^i(\bar X,\mathcal O_{\bar X})^{1-\mathfrak F_p}
\otimes_{\mathbb F_p}\mathbb F_q,
\end{equation}

el automorfismo $\varphi_q$ del miembro izquierdo de (2.2.14) se
convierte, por medio de (2.2.14), en el automorfismo inducido por
$\varphi_q$ sobre
$H^i(\bar X,\mathcal O_{\bar X})^{1-\mathfrak F_p}
\otimes_{\mathbb F_p}\mathbb F_q$. En consecuencia, el automorfismo
inducido por $\mathfrak F_q$ sobre
$H^i(\bar X,\mathcal O_{\bar X})_{\mathrm{ss}}$ se convierte, por medio
de (2.1.1), en el automorfismo inducido por
$\sisi{\varphi_q}{\varphi_q^{-1}}$ sobre
$H^i_{\mathrm{et}}(\bar X,\mathbb Z/p\mathbb Z)
\otimes_{\mathbb F_p}k$. Esto demuestra la Proposición~2.2.5.

\pageoriginale

\section*[3. Fórmula de congruencia]{3. La fórmula de congruencia: enunciados y equivalencias elementales}

Deseo agradecer a Grothendieck y Deligne que me explicaran las variantes
\sisi{(6.1.2), (3.1.1.2) y (3.1.3)}
{(3.1.2), (3.2.1)--(3.2.2) y (3.3.1)} de (3.1.1).

\subsection*{3.0.}
Sea $X$ un esquema de tipo finito sobre $\mathbb F_q$. El grado de un
punto cerrado $y\in X$ se define como el grado de su cuerpo residual
$\mathbb F_q(y)/\mathbb F_q$. La función zeta de $X/\mathbb F_q$,
considerada como elemento de $\mathbb Z[[t]]$, se define por la fórmula
(cf.~también SGA~5, XIII~3.1)

\begin{equation}\tag{3.0.1}
Z(t,X/\mathbb F_q)=\prod_y(1-t^{\deg(y)})^{-1}.
\end{equation}

Por «función zeta módulo $p$» entenderemos siempre la imagen de
$Z(t,X/\mathbb F_q)$ en $\mathbb F_q[[t]]$ por la aplicación de
«reducción módulo $p$»

\[
\mathbb Z[[t]]\longrightarrow\mathbb F_q[[t]].
\]

\subsubsection*{Teorema 3.1.}
Sea $X$ un esquema propio sobre $\mathbb F_q$. Entonces

\begin{equation}\tag{3.1.1}
Z(t,X/\mathbb F_q)
\equiv
\prod_{i\geq0}
\det\!\left(1-t\mathfrak F_q\mid H^i(X,\mathcal O_X)\right)^{(-1)^{i+1}}
\pmod p,
\end{equation}

o, equivalentemente por la Proposición~2.2.4,

\begin{equation}\tag{3.1.2}
Z(t,X/\mathbb F_q)
\equiv
\prod_{i\geq0}
\det\!\left(1-t\varphi_q^{-1}\mid
H^i_{\mathrm{et}}(\bar X,\mathbb Z/p\mathbb Z)\right)^{(-1)^{i+1}}
\pmod p.
\end{equation}

\subsubsection*{Corolario 3.2.}
Sea $X$ un esquema propio sobre $\mathbb F_q$. Entonces

\begin{equation}\tag{3.2.1}
\sum_{i\geq0}(-1)^i
\operatorname{Tr}\!\left(\mathfrak F_q\mid H^i(X,\mathcal O_X)\right)
\equiv \#X(\mathbb F_q)\pmod p,
\end{equation}

o, equivalentemente por la Proposición~2.2.4,

\pageoriginale

\begin{equation}\tag{3.2.2}
\sum_{i\geq0}(-1)^i
\operatorname{Tr}\!\left(\varphi_q^{-1}\mid
H^i_{\mathrm{et}}(\bar X,\mathbb Z/p\mathbb Z)\right)
\equiv \#X(\mathbb F_q)\pmod p.
\end{equation}

En efecto, basta tomar los coeficientes de $t$ en las congruencias (3.1.1)
y (3.1.2).

\subsubsection*{Teorema 3.3.}
Sea $X$ un esquema propio sobre $\mathbb F_q$, y denotemos por
$H^i_{\mathrm{DR}}(X/\mathbb F_q)$ los espacios vectoriales sobre
$\mathbb F_q$ $H^i(X,\Omega^\bullet_{X/\mathbb F_q})$, sobre los cuales
actúa $\mathfrak F_q$. Entonces

\begin{equation}\tag{3.3.1}
Z(t,X/\mathbb F_q)
\equiv
\prod_{i\geq0}
\det\!\left(1-t\mathfrak F_q\mid
H^i_{\mathrm{DR}}(X/\mathbb F_q)\right)^{(-1)^{i+1}}
\pmod p.
\end{equation}

En efecto, considérese la sucesión espectral

\begin{equation}\tag{3.3.2}
E_1^{p,q}=H^q(X,\Omega^p_{X/\mathbb F_q})
\Longrightarrow H^{p+q}_{\mathrm{DR}}(X/\mathbb F_q).
\end{equation}

El operador $\mathfrak F_q$ actúa sobre ella mediante el endomorfismo
$\mathfrak F_q=(\sisi{\mathrm{Fr}_p^X}{\mathrm{Fr}_q^X})^*$ inducido por
$\mathrm{Fr}_q^X:X\longrightarrow X$, y anula $E_1^{p,q}$ para
\sisi{$q\ne0$}{$p\ne0$}. Se sigue que (3.3.1) equivale a (3.1.1).

Añadamos un pequeño complemento. Como la formación de la «parte
semisimple» (cf.~1.0.7) es un funtor exacto, (3.3.2) da una sucesión
espectral

\begin{equation}\tag{3.3.3}
(E_1^{p,q})_{\mathrm{ss}}
\Longrightarrow H^{p+q}_{\mathrm{DR}}(X/\mathbb F_q)_{\mathrm{ss}}.
\end{equation}

Como $(E_1^{p,q})_{\mathrm{ss}}=0$ para
\sisi{$q\ne0$}{$p\ne0$}, obtenemos para todo $i$ un isomorfismo

\begin{equation}\tag{3.3.4}
H^i_{\mathrm{DR}}(X/\mathbb F_q)_{\mathrm{ss}}
\xrightarrow{\ \sim\ }
H^i(X,\mathcal O_X)_{\mathrm{ss}}
\end{equation}

que conmuta con $\mathfrak F_q$. Por tanto,

\begin{equation}\tag{3.3.5}
\det\!\left(1-t\mathfrak F_q\mid H^i_{\mathrm{DR}}(X/\mathbb F_q)\right)
=
\det\!\left(1-t\mathfrak F_q\mid H^i(X,\mathcal O_X)\right).
\end{equation}

\pageoriginale

\subsubsection*{Teorema 3.4.}
Sea $X$ un esquema separado de tipo finito sobre $\mathbb F_q$. Denotando
por $H_c^i$ la cohomología étale con soportes propios (SGA~4,
XVII~5.1.9(iii)), se tiene

\begin{equation}\tag{3.4.1}
Z(t,X/\mathbb F_q)
\equiv
\prod_{i\geq0}
\det\!\left(1-t\varphi_q^{-1}\mid
H_c^i(\bar X,\mathbb Z/p\mathbb Z)\right)^{(-1)^{i+1}}
\pmod p.
\end{equation}

Comencemos por demostrar la equivalencia de (3.4.1) y (3.1.1). Por el
momento, escribamos

\begin{equation}\tag{3.4.2}
Z(X)\quad\text{en lugar de }Z(t,X/\mathbb F_q)
\end{equation}

y

\begin{equation}\tag{3.4.3}
Z'(X)\quad\text{en lugar de}\quad
\prod_i\det\!\left(1-t\varphi_q^{-1}\mid
H_c^i(\bar X,\mathbb Z/p\mathbb Z)\right)^{(-1)^{i+1}}.
\end{equation}

\subsubsection*{Lema 3.5.}
Si $\{U_i\}$ es un recubrimiento finito de $X$ por subesquemas abiertos,
entonces

\begin{equation}\tag{3.5.1}
Z'(X)=
\prod_{j\geq0}
\prod_{1\leq i_0<\cdots<i_j\leq n}
Z'(U_{i_0}\cap\cdots\cap U_{i_j})^{(-1)^j}.
\end{equation}

\medskip
\noindent\emph{Demostración.}
Esto se deduce inmediatamente de la sucesión espectral de localización
(SGA~4, XVII~6.2.10)

\[
E_1^{-a,b}=\bigoplus_{i_0<\cdots<i_b}
H_c^a(\bar U_{i_0}\cap\cdots\cap\bar U_{i_b},\mathbb Z/p\mathbb Z)
\Longrightarrow H_c^{-a+b}(\bar X,\mathbb Z/p\mathbb Z).
\]

También puede obtenerse por inducción sobre $n$, usando el Lema~3.6.

\subsubsection*{Lema 3.6.}
Si $U$ es un subesquema abierto de $X$ e $Y$ es el subesquema cerrado
$X-U$, entonces

\begin{equation}\tag{3.6.1}
Z'(X)=Z'(U)Z'(Y).
\end{equation}

\medskip
\noindent\emph{Demostración.}
Esto se deduce de la sucesión exacta de cohomología

\begin{equation}\tag{3.6.2}
\longrightarrow H_c^i(\bar U,\mathbb Z/p\mathbb Z)
\longrightarrow H_c^i(\bar X,\mathbb Z/p\mathbb Z)
\longrightarrow H_c^i(\bar Y,\mathbb Z/p\mathbb Z)\longrightarrow.
\end{equation}

\pageoriginale

\subsubsection*{Lema 3.7 (Möbius).}
Sea $B$ un conjunto finito y sea $\{A_i\}$ un recubrimiento finito por
subconjuntos. Entonces

\begin{equation}\tag{3.7.1}
\operatorname{card}(B)=
\sum_{j\geq0}(-1)^j
\sum_{i_0<\cdots<i_j}
\operatorname{card}(A_{i_0}\cap\cdots\cap A_{i_j}).
\end{equation}

\medskip
\noindent\emph{Demostración.}
Hay una sucesión espectral de Leray para la cohomología racional del
espacio discreto $B$ y del recubrimiento abierto $\{A_i\}$,

\[
E_1^{p,q}=\bigoplus_{i_0<\cdots<i_q}
H^p(A_{i_0}\cap\cdots\cap A_{i_q},\mathbb Q)
\Longrightarrow H^{p+q}(B,\mathbb Q),
\]

y se calcula la característica de Euler--Poincaré de $B$.

\subsubsection*{Lema 3.8.}
Si $\{U_i\}$ es un recubrimiento finito de $X$ por subesquemas, entonces

\begin{equation}\tag{3.8.1}
Z(X)=
\prod_{j\geq0}
\prod_{i_0<\cdots<i_j}
Z(U_{i_0}\cap\cdots\cap U_{i_j})^{(-1)^j}.
\end{equation}

\medskip
\noindent\emph{Demostración.}
Se tiene

\begin{equation}\tag{3.8.2}
\operatorname{card}(X(\mathbb F_{q^n}))
=\sum_{r\mid n}r\,\operatorname{card}\{y\in X\mid\deg(y)=r\},
\end{equation}

lo cual da

\begin{equation}\tag{3.8.3}
Z(X)=\exp\!\left(
\sum_{n\geq1}\operatorname{card}(X(\mathbb F_{q^n}))\frac{t^n}{n}
\right).
\end{equation}

La conclusión se deduce de (3.7.1).

\subsubsection*{Lema 3.9.}
Supongamos que la fórmula (3.4.1), o equivalentemente (3.1.2) y por tanto
(3.1.1), vale cuando $X$ es una hipersuperficie proyectiva. Entonces vale
para todo esquema separado $X$ de tipo finito.

\pageoriginale

\medskip
\noindent\emph{Demostración.}
Por los Lemas~3.8 y 3.6, (3.4.1) vale para el complemento
$\mathbb P^n-H$ de una hipersuperficie. Como

\[
\bigcap_i(\mathbb P^n-H_i)=\mathbb P^n-\bigcup_iH_i
\]

es de nuevo el complemento de una hipersuperficie, el recubrimiento
abierto $\{\mathbb P^n-H_i\}$ de $\mathbb P^n-\bigcap_iH_i$ permite
aplicar los Lemas~3.8 y 3.5 y deducir \sisi{(3.1.3)}{(3.4.1)} para
$\mathbb P^n-\bigcap_iH_i$. Por los Lemas~3.8 y 3.6, (3.4.1) vale entonces
para $\bigcap_iH_i$, y por tanto para todo $X$ proyectivo. Otra aplicación
de los Lemas~3.8 y 3.5 da \sisi{(3.1.3)}{(3.4.1)} para todo $X$ cuasiproyectivo, en
particular para todo $X$ afín. Por último, si $X$ es separado y de tipo
finito, elíjase un recubrimiento finito por subesquemas abiertos afines y
aplíquense los Lemas~3.8 y 3.5.

\section*[4. Función zeta: método de Dwork]{4. Cálculo de la función zeta de una hipersuperficie por el
método de Dwork}

\subsection*{4.0.}
Sea

\begin{equation}\tag{4.0.1}
f(X)=f(X_1,\ldots,X_{n+2})\in\mathbb F_q[X]
\end{equation}

homogéneo de grado $d$, y supongamos que su lugar de ceros define la
hipersuperficie $V\subset\mathbb P^{n+1}$. Denotemos por $V^*$ el abierto
de $V$ en el que $\prod_{i=1}^{n+2}X_i$ es invertible, y por
$V^*_{\mathrm{aff}}$ el cono afín sobre $V^*$.

Elijamos un carácter aditivo no trivial

\begin{equation}\tag{4.0.2}
\chi:\mathbb F_{q^\nu}\longrightarrow\mu_p(C),
\end{equation}

donde $C$ es un cuerpo algebraicamente cerrado fijo de característica cero
(por ejemplo, $C=\mathbb C$). Así,

\begin{equation}\tag{4.0.3}
\text{para }x\in\mathbb F_{q^\nu},\qquad
\sum_{y\in\mathbb F_{q^\nu}}\chi(xy)=
\begin{cases}
q^\nu,&x=0,\\
0,&x\ne0.
\end{cases}
\end{equation}

Para $x=(x_1,\ldots,x_{n+2})\in(\mathbb F_{q^\nu}^*)^{n+2}$ tenemos,
por tanto,

\pageoriginale

\begin{equation}\tag{4.0.4}
1+\sum_{x_0\in\mathbb F_{q^\nu}^*}\chi(x_0f(x))
=\sum_{x_0\in\mathbb F_{q^\nu}}\chi(x_0f(x))
=
\begin{cases}
q^\nu,&f(x)=0,\\
0,&\text{en otro caso}.
\end{cases}
\end{equation}

Sumando (4.0.4) sobre $x\in(\mathbb F_{q^\nu}^*)^{n+2}$, se obtiene

\begin{equation}\tag{4.0.5}
\operatorname{card}V^*_{\mathrm{aff}}(\mathbb F_{q^\nu})
=\frac1{q^\nu}
\sum_{x\in(\mathbb F_{q^\nu}^*)^{n+2}}
\left\{1+\sum_{x_0\in\mathbb F_{q^\nu}^*}\chi(x_0f(x))\right\}.
\end{equation}

Como

\begin{equation}\tag{4.0.6}
\operatorname{card}V^*_{\mathrm{aff}}(\mathbb F_{q^\nu})
=(q^\nu-1)\operatorname{card}V^*(\mathbb F_{q^\nu}),
\end{equation}

la fórmula (4.0.5) se convierte en

\begin{equation}\tag{4.0.7}
(q^\nu-1)\operatorname{card}V^*(\mathbb F_{q^\nu})
=\frac{(q^\nu-1)^{n+2}}{q^\nu}
+\frac1{q^\nu}
\sum_{(x_0,x)\in(\mathbb F_{q^\nu}^*)^{n+3}}\chi(x_0f(x)).
\end{equation}

Escribamos explícitamente $f(X)$ como una suma de monomios,

\begin{equation}\tag{4.0.8}
f(X)=\sum_w a_wX^w,
\end{equation}

donde $w=(w_1,\ldots,w_{n+2})$, $\sum_iw_i=d$, y
$X^w\mathrel{\mathop:}=X_1^{w_1}\cdots X_{n+2}^{w_{n+2}}$. Entonces
(4.0.7) puede reescribirse como

\begin{equation}\tag{4.0.9}
(q^\nu-1)\operatorname{card}V^*(\mathbb F_{q^\nu})
=\frac{(q^\nu-1)^{n+2}}{q^\nu}
+\frac1{q^\nu}
\sum_{(x_0,x)\in(\mathbb F_{q^\nu}^*)^{n+3}}
\prod_w\chi(a_wx_0X^w).
\end{equation}

\subsection*{4.1.}
Ahora debemos explicitar nuestro «carácter aditivo con valores en
característica cero». Querríamos tomar simplemente

\begin{equation}\tag{4.1.0}
\chi:x\longmapsto
\exp\!\left(2\pi i\,\operatorname{tr}_{\mathbb F_{q^\nu}/\mathbb F_p}(x)\right)
=\xi_p^{\operatorname{tr}_{\mathbb F_{q^\nu}/\mathbb F_p}(x)}.
\end{equation}

\pageoriginale

Sea

\begin{align}
\tag{4.1.1}
\Omega&=\text{una clausura algebraica de }\mathbb Q_p,\\
\tag{4.1.2}
\operatorname{ord}:\Omega^*&\longrightarrow\mathbb Q
  &&\text{la valuación normalizada por }\operatorname{ord}(p)=1,\\
\tag{4.1.3}
\xi_p&\in\Omega
  &&\text{una raíz $p$-ésima primitiva de la unidad.}
\end{align}

Se comprueba (cf.~[2]) que existe

\begin{equation}\tag{4.1.4}
\pi\in\mathbb Q_p(\xi_p),\qquad
\operatorname{ord}(\pi)=\frac1{p-1},
\end{equation}

tal que

\begin{equation}\tag{4.1.5}
\sum_{n\geq0}\frac{\pi^{p^n}}{p^n}=0.
\end{equation}

(Obsérvese también que el miembro de la izquierda converge en
$\mathbb Q_p(\xi_p)$.)

Definamos la exponencial de Artin--Hasse

\begin{equation}\tag{4.1.6}
E(X)=\exp\!\left(\sum_{n\geq0}\frac{X^{p^n}}{p^n}\right)
\in\mathbb Q_p[[X]],
\end{equation}

que, milagrosamente, pertenece a $\mathbb Z_p[[X]]$. Definamos

\begin{equation}\tag{4.1.7}
\Theta(X)=E(\pi X).
\end{equation}

Sus propiedades fundamentales se resumen en el lema siguiente.

\subsubsection*{Lema 4.2.}

\begin{equation}\tag{4.2.1}
\Theta(X)=\sum_n B_n X^n,\qquad B_0=1,\quad
\sisi{A_1}{B_1}=\pi,\quad
\operatorname{ord}(B_n)\geq\frac n{p-1}.
\end{equation}

\begin{equation}\tag{4.2.2}
\Theta(1)=\xi_p,\qquad\text{una raíz $p$-ésima primitiva de la unidad.}
\end{equation}

Sea $t\in\Omega$ y supongamos que, para cierto entero $a>0$,

\begin{equation}\tag{4.2.3}
t=t^{p^a};
\end{equation}

es decir, $t$ es un representante de Teichmüller en
$W(\mathbb F_{\sisi{q^a}{p^a}})\subset\Omega$. Entonces
$\sum_{j=0}^{a-1}t^{p^j}\in\mathbb Z_p$, de modo que la serie formal

\pageoriginale

\[
(\Theta(X))^{\sum_{j=0}^{a-1}\sisi{t^j}{t^{p^j}}}
\]

está definida, y se tiene la identidad de series formales

\begin{equation}\tag{4.2.3.1}
\prod_{j=0}^{a-1}\Theta(t^{p^j}X)
=(\Theta(X))^{\sum_{j=0}^{a-1}\sisi{t^j}{t^{p^j}}}.
\end{equation}

\medskip
\noindent\emph{Demostración.}
Como $E(X)\in\mathbb Z_p[[X]]$, (4.2.1) se deduce de la definición. Para
(4.2.2), se tiene

\begin{equation}\tag{4.2.4}
\log(\Theta(1))=0
\qquad\text{por la elección de $\pi$ en (4.1.5),}
\end{equation}

y

\begin{equation}\tag{4.2.5}
\operatorname{ord}(\Theta(1)-1)=\frac1{p-1}
\qquad\text{por (4.2.1).}
\end{equation}

Por (4.2.4), existe un entero $n\geq0$ tal que
$\Theta(1)^{p^n}=1$ (esta es la caracterización del núcleo del logaritmo),
y (4.2.5) implica que $n=1$.

En cuanto a (4.2.3.1), tras tomar logaritmos se deduce por un cálculo
inmediato, una vez que se sabe que
$\sum_{j=0}^{a-1}t^{p^j}\in\mathbb Z_p$. Pero los $t^{p^j}$,
$0\leq j\leq a-1$, son precisamente los conjugados sobre $\mathbb Z_p$
de $t\in W(\mathbb F_{p^a})$; por tanto,

\[
\sum_{j=0}^{a-1}t^{p^j}
=\operatorname{Tr}_{W(\mathbb F_{p^a})/\mathbb Z_p}(t)\in\mathbb Z_p.
\]

Evaluando en $X=1$ ambos miembros de (4.2.3.1), se obtiene el resultado
siguiente.

\subsubsection*{Corolario 4.3.}
Sea $t\in\Omega$ tal que $t^{p^a}=t$, y denotemos su clase residual por
$\bar t\in\mathbb F_{p^a}$. Entonces

\begin{equation}\tag{4.3.1}
\prod_{j=0}^{a-1}\Theta(t^{p^j})
=\xi_p^{\operatorname{tr}_{\mathbb F_{p^a}/\mathbb F_p}(\bar t)}.
\end{equation}

\pageoriginale

\subsection*{4.4.}
Recordemos que el cuerpo residual del anillo de enteros de
$\mathbb Q_p(\xi_p)$ es $\mathbb F_p$. Sea

\begin{equation}\tag{4.4.0}
K=\text{la extensión no ramificada de $\mathbb Q_p(\xi_p)$ con cuerpo
residual $\mathbb F_q$},
\end{equation}

y sea

\begin{equation}\tag{4.4.1}
\tau\in\operatorname{Gal}(K/\mathbb Q_p(\xi_p))
\end{equation}

el «Frobenius», es decir, el automorfismo cuya reducción es el automorfismo
de Frobenius de $\mathbb F_q/\mathbb F_p$, $x\longmapsto x^p$.

Sea

\begin{equation}\tag{4.4.2}
F(X)=\sum_w A_wX^w\in K[X_1,\ldots,X_{n+2}]
\end{equation}

el levantamiento de $f(X)$ cuyos coeficientes satisfacen $A_w^q=A_w$.
Reuniendo (4.3.1) y (4.0.9), y teniendo en cuenta la definición (4.1.0) de
$\chi$, se obtiene

\begin{equation}\tag{4.4.3}
\begin{aligned}
&(q^\nu-1)\operatorname{card}V^*(\mathbb F_{q^\nu})
-\frac{(q^\nu-1)^{n+2}}{q^\nu}\\
&\quad=\frac1{q^\nu}
\sum_{(X_0,X)\in(\mu_{q^\nu-1}(\Omega))^{n+3}}
\prod_w\prod_{j=0}^{\operatorname{ord}(q^\nu)-1}
\Theta\!\left(A_w^{p^j}X_0^{p^j}X^{p^jw}\right).
\end{aligned}
\end{equation}

Definamos

\[
H(X)=H(X_0,\ldots,X_{n+2})\in K[[X_0,\ldots,X_{n+2}]]
\]

por

\begin{equation}\tag{4.4.4}
H(X)=\prod_w
\Theta\!\left(A_wX_0X_1^{w_1}\cdots X_{n+2}^{w_{n+2}}\right).
\end{equation}

Hacemos actuar $\operatorname{Gal}(K/\mathbb Q_p(\xi_p))$ sobre
$K[[X_0,\ldots,X_{n+2}]]$ mediante su acción sobre los coeficientes; la
acción componente a componente de
$\sigma\in\operatorname{Gal}(K/\mathbb Q_p(\xi_p))$ se denotará por

\begin{equation}\tag{4.4.5}
G(X)\longmapsto G^\sigma(X).
\end{equation}

\pageoriginale

En particular,

\[
H^{\tau^j}(X)=\prod_w
\Theta\!\left(A_w^{\tau^j}X_0X_1^{w_1}\cdots
X_{n+2}^{w_{n+2}}\right).
\]

Como $A_w^{\tau^j}=A_w^{p^j}$, la fórmula (4.4.3) puede reescribirse
como

\begin{equation}\tag{4.4.6}
\begin{aligned}
(q^\nu-1)\operatorname{card}V^*(\mathbb F_{q^\nu})
&=\frac{(q^\nu-1)^{n+2}}{q^\nu}\\
&\quad+\frac1{q^\nu}
\sum_{(X_0,X)\in(\mu_{q^\nu-1}(\Omega))^{n+3}}
\prod_{j=0}^{\operatorname{ord}(q^\nu)-1}H^{\tau^j}(X^{p^j}).
\end{aligned}
\end{equation}

\subsection*{4.5.}
Ahora expresaremos la suma de (4.4.6) como una traza.

Para cada $b\in\mathbb Q$, definamos un álgebra de Banach $p$-ádica sobre
$K$

\begin{equation}\tag{4.5.1}
L(b)=\text{el subanillo de $K[[X_0,\ldots,X_{n+2}]]$ formado por las series}
\end{equation}

que pueden escribirse

\begin{equation}\tag{4.5.2}
\sum_{dU_0=\sum_{i\geq1}U_i}
B_UX_0^{U_0}X_1^{U_1}\cdots X_{n+2}^{U_{n+2}},
\qquad U=(U_0,\ldots,U_{n+2}),
\end{equation}

y que satisfacen

\begin{equation}\tag{4.5.3}
\inf_U\bigl(\operatorname{ord}(B_U)-bU_0\bigr)>-\infty.
\end{equation}

Es un álgebra de Banach $p$-ádica para la función orden (el opuesto del
logaritmo de la norma)

\begin{equation}\tag{4.5.4}
\inf_U\bigl(\operatorname{ord}(B_U)-bU_0\bigr)
=\operatorname{ord}\!\left(\sum_U B_UX^U\right).
\end{equation}

Se deduce de (4.4.4) y (4.2.1) que

\[
H(X)\in L\!\left(\frac1{p-1}\right).
\]

\pageoriginale

Introduzcamos algunos operadores. Para todo entero $k>0$, definamos un
operador continuo entre espacios de Banach sobre $K$ por

\begin{equation}\tag{4.5.5}
\left\{
\begin{aligned}
\psi_k:L(b)&\longrightarrow L(kb),\\
\psi_k\!\left(\sum_U B_UX^U\right)&=\sum_U B_{kU}X^U.
\end{aligned}
\right.
\end{equation}

Para todo par $(b,b')$ con $b'<b$, hay un operador de «restricción», que
es completamente continuo (es decir, límite uniforme de operadores de
rango finito):

\begin{equation}\tag{4.5.6}
\left\{
\begin{aligned}
L(b)&\xrightarrow{\ \text{restricción}\ }L(b'),\\
\sum_U B_UX^U&\longmapsto\sum_U B_UX^U.
\end{aligned}
\right.
\end{equation}

Cada elemento $G$ del álgebra $L(b)$ da el endomorfismo continuo de
multiplicación por $G$,

\begin{equation}\tag{4.5.7}
\left\{
\begin{aligned}
G:L(b)&\longrightarrow L(b),\\
\eta&\longmapsto\eta G.
\end{aligned}
\right.
\end{equation}

Estos operadores están relacionados por el lema siguiente, cuya
demostración se encuentra en [2] (véase también [4]).

\pageoriginale

\subsubsection*{Lema 4.5.8 (lema de la traza de Dwork).}
Sean $b\in\mathbb Q$, $b>0$, $k\in\mathbb Z$, $k\geq2$, y
$G\in L(b/k)$. Entonces el compuesto

\[
L(b)\xrightarrow{\ \text{restricción}\ }L(b/k)
\xrightarrow{\ \times G\ }L(b/k)
\xrightarrow{\ \psi_k\ }L(b),
\]

denotado por $\psi_k\circ G$, es completamente continuo, y su traza
(cf.~[4]) viene dada por

\begin{equation}\tag{4.5.9}
\operatorname{tr}(\psi_k\circ G)
=\frac1{(k-1)^{n+2}}
\sum_{X\in(\mu_{k-1}(\Omega))^{n+3}}G(X).
\end{equation}

Apliquemos este lema con

\[
\left\{
\begin{aligned}
k&=q^\nu,\\
G(X)&=\prod_{j=0}^{\operatorname{ord}(q^\nu)-1}
H^{\tau^j}(X^{p^j})
\in L\!\left(\frac{p}{(p-1)q^\nu}\right),\\
b&=\frac1{p-1}.
\end{aligned}
\right.
\]

La fórmula (4.4.6) se convierte en

\begin{equation}\tag{4.5.10}
\begin{aligned}
(q^\nu-1)\operatorname{card}V^*(\mathbb F_{q^\nu})
&=\frac{(q^\nu-1)^{n+2}}{q^\nu}\\
&\quad+\frac{(\sisi{q}{q^\nu}-1)^{n+3}}{q^\nu}
\operatorname{tr}\!\left(
\left.\psi_{q^\nu}\circ
\prod_{j=0}^{\operatorname{ord}(q^\nu)-1}H^{\tau^j}(X^{p^j})
\right|L\!\left(\frac1{p-1}\right)\right).
\end{aligned}
\end{equation}

Sea $L^0(b)$ el ideal de los elementos de $L(b)$ que «se anulan en
cero», es decir, que no tienen término constante. Es estable por
$\psi_k\circ G$, y evidentemente

\begin{equation}\tag{4.5.11}
\operatorname{tr}(\psi_k\circ G)
=G(0)+\operatorname{tr}\!\left(\left.\psi_k\circ G\right|L^0(b)\right).
\end{equation}

Como $H(0)=1$, la fórmula (4.5.10) puede reescribirse como

\begin{equation}\tag{4.5.12}
\begin{aligned}
\operatorname{card}V^*(\mathbb F_{q^\nu})
&=(q^\nu-1)^{n+1}\\
&\quad+\frac{(q^\nu-1)^{n+2}}{\sisi{q}{q^\nu}}
\operatorname{tr}\!\left(
\left.\psi_{q^\nu}\circ
\prod_{j=0}^{\operatorname{ord}(q^\nu)-1}H^{\tau^j}(X^{p^j})
\right|L^0\!\left(\frac1{p-1}\right)\right).
\end{aligned}
\end{equation}

\pageoriginale

Consideremos ahora el operador $\tau^{-1}$-lineal

\begin{equation}\tag{4.5.13}
\tau^{-1}\circ\psi_p\circ H(X):L(b)\longrightarrow L(b).
\end{equation}

Sus iteradas son

\begin{equation}\tag{\sisi{5.4.14}{4.5.14}}
(\tau^{-1}\circ\psi_p\circ H(X))^a
=\tau^{-a}\circ\psi_{p^a}\circ
H(X)H^\tau(X^p)\cdots H^{\tau^{a-1}}(X^{p^{a-1}}),
\end{equation}

y, si $\operatorname{ord}(q)$ divide a $a$, entonces
$\tau^{-a}=\operatorname{id}$. Definamos la operación $\tau^{-1}$-lineal

\begin{equation}\tag{4.5.15}
\left\{
\begin{aligned}
\alpha:L^0\!\left(\frac1{p-1}\right)
&\longrightarrow L^0\!\left(\frac1{p-1}\right),\\
\alpha&=\frac1p\,\tau^{-1}\circ\psi_p\circ H(X).
\end{aligned}
\right.
\end{equation}

Entonces (4.5.12) toma la forma

\begin{equation}\tag{4.5.16}
\operatorname{card}V^*(\mathbb F_{q^\nu})
=(q^\nu-1)^{n+1}
+(q^\nu-1)^{n+2}
\operatorname{tr}\!\left(
\left.\alpha^{\operatorname{ord}(q^\nu)}
\right|L^0\!\left(\frac1{p-1}\right)\right).
\end{equation}

Pongamos

\begin{equation}\tag{4.5.17}
S=\{1,\ldots,n+2\}.
\end{equation}

Para todo subconjunto no vacío $A\subset S$, definamos el subespacio

\begin{equation}\tag{4.5.18}
\begin{aligned}
L^A(b)=\{\,&\text{series $\sum_U B_UX^U\in L(b)$ tales que, para cada
monomio $X^U$ con $B_U\ne0$,}\\
&\text{los índices $i$ para los que $U_i>0$ son exactamente $0$ y los
$j\in A$}\,\}.
\end{aligned}
\end{equation}

Entonces

\begin{equation}\tag{4.5.19}
L^0(b)=\bigoplus_{\substack{A\subset S\\A\ne\varnothing}}L^A(b).
\end{equation}

\pageoriginale

Y, para todo subconjunto no vacío $B\subset S$, se tiene

\begin{equation}\tag{4.5.20}
\bigoplus_{A\supset B}L^A(b)
=L(b)\cap\left(\prod_{i\in B}X_i\right)K[[X]].
\end{equation}

Evidentemente, cada subespacio $\bigoplus_{A\supset B}L^A(b)$ es estable
por $\alpha$. Se sigue que lo mismo ocurre con
$\bigoplus_{A\not\supset B}L^A(b)$, que es una suma de subespacios del
tipo anterior. Definamos

{\small
\begin{equation}\tag{4.5.21}
\alpha_B=\text{el endomorfismo $\tau^{-1}$-lineal inducido por $\alpha$
sobre}\quad
\frac{\displaystyle\bigoplus_{A\supset B}
L^A\!\left(\frac1{p-1}\right)}
{\displaystyle\bigoplus_{\substack{A\supset B\\A\ne B}}
L^A\!\left(\frac1{p-1}\right)}
\simeq L^B\!\left(\frac1{p-1}\right).
\end{equation}
}

Explícitamente, $\alpha_B$ viene dado por

\begin{equation}\tag{4.5.22}
\alpha_B=\frac{\tau^{-1}}p\,\psi_p\,H_B(X),
\end{equation}

operación bajo la cual $L^B$ es estable, donde $H_B(X)$ se obtiene de
$H(X)$ sustituyendo $X_i=0$ para $i\notin\{0\}\cup B$.

Reescribiendo una vez más (4.5.16), se obtiene

\begin{equation}\tag{\sisi{4.5.2.3}{4.5.23}}
\operatorname{card}V^*(\mathbb F_{q^\nu})
=(q^\nu-1)^{n+1}
+(q^\nu-1)^{n+2}
\sum_{\substack{A\subset S\\A\ne\varnothing}}
\operatorname{tr}\!\left((\alpha_A)^{\operatorname{ord}(q^\nu)}\right).
\end{equation}

Para $A\subset S$, $A\ne\varnothing$, definamos la hipersuperficie

\begin{equation}\tag{4.5.24}
V_A\subset\mathbb P^{\operatorname{card}(A)-1}
\simeq
\mathbb P^{n+1}\cap
\{X_i=0\text{ para todo }i\notin A\}
\end{equation}

mediante $f_A(X)$, el polinomio obtenido de $f(X)$ sustituyendo $X_i=0$
para $i\notin A$.

\pageoriginale

Por la compatibilidad de la formación de la serie $H$ a partir de $f(X)$
con la sustitución por $0$ de algunas variables, (4.5.23) da, para todo
$A\subset S$, $A\ne\varnothing$,

\begin{equation}\tag{4.5.25}
\operatorname{card}V_A^*(\mathbb F_{q^\nu})
=(q^\nu-1)^{\sisi{\operatorname{car}(A)}{\operatorname{card}(A)}-1}
+(q^\nu-1)^{\operatorname{card}(A)}
\sum_{B\subset A}
\operatorname{tr}\!\left((\alpha_B)^{\operatorname{ord}(q^\nu)}\right).
\end{equation}

Por otra parte,

\begin{equation}\tag{4.5.26}
\operatorname{card}V(\mathbb F_{q^\nu})
=\sum_{\substack{A\subset S\\A\ne\varnothing}}
\operatorname{card}V_A^*(\mathbb F_{q^\nu}).
\end{equation}

Después de una suma sencilla, (4.5.25) da por tanto

\begin{equation}\tag{4.5.27}
\begin{aligned}
\operatorname{card}V(\mathbb F_{q^\nu})
={}&1+q^\nu+q^{2\nu}+\cdots
+q^{\nu(\operatorname{card}(S)-1)}\\
&+\sum_{\substack{B\subset S\\B\ne\varnothing}}
(q^\nu-1)^{\operatorname{card}(B)}
q^{\nu\operatorname{card}(S-B)}
\operatorname{tr}\!\left((\alpha_B)^{\operatorname{ord}(q^\nu)}\right).
\end{aligned}
\end{equation}

Como $\alpha$, y por tanto los $\alpha_B$, son $\tau^{-1}$-lineales, los
operadores $\alpha_B^{\operatorname{ord}(q)}$ son $K$-lineales. Definamos
la serie característica

\[
\Delta_B(t)\in1+tK[[t]]
\]

por

\begin{equation}\tag{4.5.28}
\Delta_B(t)=\det\!\left(
1-t\alpha_B^{\operatorname{ord}(q)}
\mathrel{\Big|}L^B\!\left(\frac1{p-1}\right)\right).
\end{equation}

Introduzcamos los endomorfismos siguientes de $1+tK[[t]]$:

\begin{equation}\tag{4.5.29}
\varphi:g(t)\longmapsto g^\varphi(t)=g(qt)
\end{equation}

y

\begin{equation}\tag{4.5.30}
\delta=1-\varphi:g(t)\longmapsto
g^\delta(t)=\frac{g(t)}{g(qt)}.
\end{equation}

\pageoriginale

Teniendo en cuenta las identidades

\begin{equation}\tag{4.5.31}
Z(t,V/\mathbb F_q)
=\exp\!\left(
\sum_{\nu\geq1}
\operatorname{card}V(\mathbb F_{q^\nu})\frac{t^\nu}{\nu}
\right)
\end{equation}

y

\begin{equation}\tag{4.5.32}
\det\!\left(1-t\alpha_B^{\operatorname{ord}(q)}\right)
=\exp\!\left(
-\sum_{\nu\geq1}
\operatorname{tr}\!\left((\alpha_B)^{\nu\operatorname{ord}(q)}\right)
\frac{t^\nu}{\nu}
\right),
\end{equation}

la fórmula (4.5.27) puede reescribirse como

\begin{equation}\tag{\sisi{4.5.23}{4.5.33}}
Z(t,V/\mathbb F_q)\prod_{i=0}^{n+1}(1-q^it)
=\prod_{j=1}^{n+2}
\left(
\prod_{\substack{B\subset S\\\operatorname{card}(B)=j}}
\Delta_B(t)
\right)^{-(-\delta)^j\varphi^{n+2-j}}.
\end{equation}

\subsection*{4.6.}
Estudiemos ahora la integralidad de las series $\Delta_B$. Pongamos

\begin{equation}\tag{4.6.1}
\mathcal O=\text{el anillo de enteros de $K$},
\end{equation}

y definamos una filtración sobre $1+t\mathcal O[[t]]$ por

\begin{equation}\tag{4.6.2}
F^i(1+t\mathcal O[[t]])
=1+q^it\mathcal O[[q^it]]
=\varphi^i(1+t\mathcal O[[t]]).
\end{equation}

Se tiene

\begin{align}
\varphi(F^i)&\subset F^{i+1}, \tag{4.6.3}\\
\delta(F^i)&\subset F^i, \tag{4.6.4}\\
\delta&\text{ induce la identidad sobre }F^i/F^{i+1}. \tag{4.6.5}
\end{align}

\subsubsection*{Lema 4.6.6.}
Se tiene $\Delta_B(t)\in1+t\mathcal O[[t]]$; además, si un entero $j$
satisface

\begin{equation}\tag{4.6.7}
jd<\operatorname{card}(B),
\end{equation}

entonces

\begin{equation}\tag{4.6.8}
\Delta_B(t)\in F^j(1+t\mathcal O[[t]]).
\end{equation}

\pageoriginale

\noindent\emph{Demostración.}
Los espacios de Banach $L^B\!\left(\frac1{p-1}\right)$ son todos del
tipo $B(I)$ en el sentido de Serre~[7]; es decir, admiten bases
ortonormales, a saber, los monomios

\begin{equation}\tag{4.6.9}
\pi^{U_0}X_0^{U_0}X_1^{U_1}\cdots X_{n+2}^{U_{n+2}}
\end{equation}

que les pertenecen. Decir que esta es una base ortonormal significa, por
definición, que

\begin{equation}\tag{4.6.10}
\operatorname{ord}\!\left(
\sum_U C_U\pi^{U_0}X_0^{U_0}\cdots X_{n+2}^{U_{n+2}}
\right)
=\min_U\operatorname{ord}(C_U).
\end{equation}

Sea

\begin{equation}\tag{4.6.11}
R_B=\text{el submódulo sobre $\mathcal O$ de
$L^B\!\left(\frac1{p-1}\right)$ formado por los elementos de orden
$\geq0$}.
\end{equation}

Para demostrar el lema basta probar

\begin{equation}\tag{4.6.12}
jd<\operatorname{card}(B)
\quad\Longrightarrow\quad
(\alpha_B)^{\operatorname{ord}(q)}R_B\subset q^jR_B,
\end{equation}

o, con mayor razón, probar

\begin{equation}\tag{4.6.13}
jd<\operatorname{card}(B)
\quad\Longrightarrow\quad
\alpha_B(R_B)\subset p^jR_B.
\end{equation}

Por (4.5.22),

\begin{equation}\tag{4.6.14}
\alpha_B=\frac1p\,\tau^{-1}\circ\psi_p\circ H_B(X).
\end{equation}

Como

\begin{equation}\tag{4.6.15}
H(X)\in L\!\left(\frac1{p-1}\right),
\qquad \operatorname{ord}(H(X))\geq0,
\end{equation}

se tiene

\begin{equation}\tag{4.6.16}
H_B(X)\,R_B\subset R_B
\end{equation}

y, evidentemente,

\begin{equation}\tag{4.6.17}
\tau^{-1}(R_B)=R_B.
\end{equation}

\pageoriginale

En consecuencia, (4.6.13) se deduce de

\begin{equation}\tag{4.6.18}
jd<\operatorname{card}(B)
\quad\Longrightarrow\quad
\psi_p(R_B)\subset p^{j+1}R_B.
\end{equation}

Verifiquemos (4.6.18) monomio a monomio. Se tiene

{\small
\begin{equation}\tag{4.6.19}
\begin{cases}
\psi_p\!\left(\pi^{U_0}X_0^{U_0}\cdots X_{n+2}^{U_{n+2}}\right)=0,
&\text{salvo si }p\mid U_i\text{ para todo }i,\\[3pt]
\psi_p\!\left(
\pi^{pU_0}X_0^{pU_0}\cdots X_{n+2}^{pU_{n+2}}
\right)
=\pi^{(p-1)U_0}
\underbrace{\pi^{U_0}X_0^{U_0}\cdots X_{n+2}^{U_{n+2}}}_{\in R_B}.
\end{cases}
\end{equation}
}

Las condiciones para que
$\pi^{U_0}X_0^{U_0}\cdots X_{n+2}^{U_{n+2}}$ pertenezca a
$L^B\!\left(\frac1{p-1}\right)$ son

\begin{equation}\tag{4.6.20}
\begin{cases}
dU_0=\displaystyle\sum_{i\geq1}U_i,\\
U_i>0\Longleftrightarrow i\in B,\qquad i\geq1.
\end{cases}
\end{equation}

En particular,

\begin{equation}\tag{4.6.21}
dU_0\geq\operatorname{card}(B).
\end{equation}

Por tanto,

\begin{equation}\tag{4.6.22}
\operatorname{ord}\!\left(\pi^{(p-1)U_0}\right)
=U_0\geq\frac{\operatorname{card}(B)}d>j,
\end{equation}

y, como $U_0\in\mathbb Z$,

\begin{equation}\tag{4.6.23}
\operatorname{ord}\!\left(\pi^{(p-1)U_0}\right)
=U_0\geq j+1.
\end{equation}

Junto con (4.6.19), esto demuestra (4.6.18).

\pageoriginale

\section*{5. Un teorema de Ax [1] (cf. también [4], sección 7)}

\subsubsection*{Teorema 5.0.}
Sea $V$ una hipersuperficie de $\mathbb P^{n+1}$ de grado $d>0$,
definida sobre $\mathbb F_q$, y sea $j$ un entero. Si

\begin{equation}\tag{5.0.1}
jd<n+2,
\end{equation}

entonces

\begin{equation}\tag{5.0.2}
Z(t,V/\mathbb F_q)\prod_{i=0}^{n+1}(1-q^it)
\in1+q^jt\mathbb Z[[q^jt]],
\end{equation}

y el inverso de todo cero o polo de

\begin{equation}\tag{5.0.3}
Z(t,V/\mathbb F_q)\prod_{i=0}^{n+1}(1-q^it)
\end{equation}

es de la forma $q^j$ por un entero algebraico.

\subsubsection*{Teorema 5.0 bis.}
Sea $V_{\mathrm{aff}}$ una hipersuperficie de $\mathbb A^{n+2}$ definida
por una forma homogénea de grado $d$, sobre $\mathbb F_q$. Si

\begin{equation}\tag{\sisi{5.0.7}{5.0.4}}
jd<n+2,
\end{equation}

entonces

\begin{equation}\tag{5.0.5}
Z(t,V_{\mathrm{aff}}/\mathbb F_q)
\in1+q^jt\mathbb Z[[q^jt]].
\end{equation}

En particular,

\begin{equation}\tag{5.0.6}
\operatorname{card}\bigl(V_{\mathrm{aff}}(\mathbb F_q)\bigr)
\equiv0\pmod{q^j},
\end{equation}

y, por consiguiente,

\begin{equation}\tag{5.0.6 bis}
\operatorname{card}\bigl(V(\mathbb F_q)\bigr)
\geq\frac{q^j-1}{q-1}>0,
\end{equation}

lo que mejora el teorema de Chevalley--Warning~[8].

\pageoriginale

\medskip
\noindent\emph{Demostración.}
La implicación (5.0.5) $\Rightarrow$ (5.0.6) es clara tomando el
coeficiente de $t$ en la función zeta. Las afirmaciones (5.0.4) y (5.0.2)
son equivalentes, pues

\begin{equation}\tag{5.0.7}
\operatorname{card}V_{\mathrm{aff}}(\mathbb F_{q^\nu})
=1+(q^\nu-1)\operatorname{card}V(\mathbb F_{q^\nu}),
\end{equation}

de donde

\begin{equation}\tag{5.0.8}
Z(t,V_{\mathrm{aff}}/\mathbb F_q)
=\frac{Z(t,V/\mathbb F_q)}
{Z(qt,V/\mathbb F_q)(1-t)}
\end{equation}

y

\begin{equation}\tag{5.0.9}
Z(t,V/\mathbb F_q)
=\prod_{i\geq0}
\left[(1-q^i\sisi{f}{t})Z(q^it,V_{\mathrm{aff}}/\mathbb F_q)\right].
\end{equation}

Las afirmaciones (5.0.2) y (5.0.3) son asimismo equivalentes, por el
teorema de Fatou~[3] y la racionalidad de la función zeta. Como

\begin{equation}\tag{5.0.10}
Z(t,V/\mathbb F_q)\in\mathbb Z[[t]],
\end{equation}

basta demostrar

\begin{equation}\tag{5.0.11}
Z(t,V/\mathbb F_q)\prod_{i=0}^{n+1}(1-q^it)
\in F^j(1+t\mathcal O[[t]]).
\end{equation}

A la vista de (4.5.33) y (4.6.4)--(4.6.5), basta mostrar, para
$\varnothing\ne B\subset\{1,\ldots,n+2\}$, que

\begin{equation}\tag{5.0.12}
\Delta_B(t)\in F^{j+\operatorname{card}(B)-(n+2)}.
\end{equation}

Podemos terminar la demostración aplicando el
Lema~\sisi{4.4.6}{4.6.6} una vez probado que

\begin{equation}\tag{5.0.13}
jd<n+2
\quad\Longrightarrow\quad
\bigl(j+\operatorname{card}(B)-(n+2)\bigr)d
<\operatorname{card}(B).
\end{equation}

Esto puede reescribirse como

\pageoriginale

\begin{equation}\tag{5.0.14}
jd<n+2
\quad\Longrightarrow\quad
(j-(n+2))d<\operatorname{card}(B)(1-d).
\end{equation}

Pero

\[
\inf_{\varnothing\ne B\subset\{1,\ldots,n+2\}}
\operatorname{card}(B)(1-d)
=(n+2)(1-d),
\]

por lo que basta demostrar

\begin{equation}\tag{5.0.15}
jd<n+2
\quad\Longrightarrow\quad
(j-(n+2))d<(n+2)(1-d),
\end{equation}

lo cual es cierto.

\subsubsection*{Observación 5.1.}
La condición

\begin{equation}\tag{5.1.1}
jd<n+2
\end{equation}

admite una interpretación cohomológica. Sea $V$ una hipersuperficie lisa
de grado $d$ y dimensión $n$. Como en XI~2.1, pongamos, para $a+b=n$,

\begin{equation}\tag{5.1.2}
h_0^{a,b}
=\dim H^b(V,\Omega_V^a)-\delta_{a,b}
=\dim E^{a,b}(V),
\end{equation}

donde $E^{a,b}(V)$ es la componente de bigrado $(a,b)$ de la cohomología
de Hodge evanescente (compárese XVII~5.4.4)

\begin{equation}\tag{5.1.3}
E^{a,b}(V)
=\operatorname{Coker}\!\left(
H^b(\mathbb P^{n+1},\Omega_{\mathbb P^{n+1}}^a)
\longrightarrow H^b(V,\Omega_V^a)\right).
\end{equation}

Entonces, por XI~2.8,

\begin{equation}\tag{5.1.4}
jd<n+2
\quad\Longleftrightarrow\quad
h_0^{a,n-a}=0\quad\text{para }a<j\text{ y }a>n-j.
\end{equation}

En otras palabras, la condición dice que el nivel de Hodge (XI~2.7) de la
cohomología de De Rham evanescente de $V$,

\[
E_{\mathrm{DR}}^n(V)
=\operatorname{Coker}\!\left(
H_{\mathrm{Hdg}}^n(\mathbb P^{n+1})
\longrightarrow H_{\mathrm{Hdg}}^n(V)\right),
\]

\pageoriginale

es a lo sumo $n-2j$, es decir, que su conivel es al menos $j$. Por otra
parte, la conclusión dice también que los valores propios del Frobenius
$\mathfrak F_q$ que actúa sobre la cohomología $\ell$-ádica evanescente
($\ell\ne p$) son enteros algebraicos divisibles por $q^j$ (XXI~A.2).
Así, el resultado 5.0 de Ax apoya la conjetura siguiente, sugerida por el
yoga del nivel de los motivos y que relaciona las conjeturas clásicas de
Hodge y de Tate~[5].

\subsubsection*{Conjetura 5.2.}
Sea $f:X\to S$ un morfismo propio y liso, con $S$ integral, de tipo finito
sobre $\operatorname{Spec}(\mathbb Z)$ y dominante sobre
$\operatorname{Spec}(\mathbb Z)$; sean $i,j$ enteros. Supongamos que el
conivel de Hodge de $H^i$ de la fibra genérica de $f$ es al menos $j$.
Entonces, para todo punto $s\in S$, el conivel (XX~2.0.6) de
$H^i(X_s,\mathbb Q_\ell)$ es al menos $j$. Por tanto (XX~2.2), si $s$ es
un punto cerrado de $S$, los valores propios de Frobenius sobre
$H^i(X_s,\mathbb Q_\ell)$ son enteros algebraicos divisibles por $q^j$,
donde $q=\operatorname{card}(k(s))$.

\section*{6. Fin de la demostración de la fórmula de congruencia}

Teniendo en cuenta (4.6.6) y (4.6.3)--(4.6.4), la fórmula
\sisi{(4.5.24)}{(4.5.33)} da

\begin{equation}\tag{6.0.1}
Z(t,V/\mathbb F_q)(1-t)
\equiv\Delta_S(t)^{(-1)^{n+1}}
\pmod{1+qt\mathcal O[[t]]}.
\end{equation}

Calculemos ahora $\Delta_S(t)$ módulo $\pi\mathcal O[[t]]$. Definamos,
como en (4.6.11),

\begin{equation}\tag{6.0.2}
W_S=\sisi{R_B}{R_S}\otimes_{\mathcal O}\mathbb F_q.
\end{equation}

\pageoriginale

\begin{equation}\tag{6.0.3}
\bar\alpha_S:W_S\longrightarrow W_S
\end{equation}

denota la operación $\tau^{-1}$-lineal inducida por $\alpha_S$;

\begin{equation}\tag{6.0.4}
R_S(i)
=\text{el submódulo de $R_S$ formado por los elementos de grado $i$ en
$X_0$},
\end{equation}

y

\begin{equation}\tag{6.0.5}
W_S(i)=R_S(i)\otimes_{\mathcal O}\mathbb F_q.
\end{equation}

\subsubsection*{Proposición 6.0.6.}
Se tiene
\[
\alpha_S(R_S)\subset R_S(1)+pR_S.
\]

\noindent\emph{Demostración.}
Se tiene la factorización (4.6.14)

\begin{equation}\tag{6.0.7}
\alpha_S=\frac1p\,\tau^{-1}\circ\psi_p\circ H(X).
\end{equation}

Sabemos (véase (4.6.15)) que

\begin{equation}\tag{6.0.8}
H(X)\,R_S\subset R_S,
\end{equation}

y, evidentemente,

\begin{equation}\tag{6.0.9}
\tau^{-1}\bigl(R_S(1)+pR_S\bigr)=R_S(1)+pR_S.
\end{equation}

Basta, por tanto, verificar

\begin{equation}\tag{6.0.10}
\frac1p\,\psi_p(R_S)\subset R_S(1)+pR_S,
\end{equation}

lo cual es claro por (4.6.19).

\subsubsection*{Corolario 6.0.11.}
Se tiene
\[
\bar\alpha_S(W_S)\subset W_S(1).
\]

\pageoriginale

\subsubsection*{Corolario 6.0.12.}
Se tiene la congruencia

\[
\Delta_S(t)\equiv
\det\!\left(
1-t(\bar\alpha_S)^{\operatorname{ord}(q)}
\mathrel{\Big|}W_S(1)\right)
\pmod{\pi\mathcal O[[t]]}.
\]

\subsubsection*{Proposición 6.1.}
La operación $\tau^{-1}$-lineal
$\bar\alpha_S:W_S(1)\to W_S(1)$ es igual a

\begin{equation}\tag{6.1.1}
\frac1p\,\tau^{-1}\circ\psi_p\circ(X_0f(X))^{p-1}.
\end{equation}

\noindent\emph{Demostración.}
Escribamos

\begin{equation}\tag{6.1.2}
H(X)=1+\sum_{i\geq1}H_i(X),
\end{equation}

donde

\begin{equation}\tag{6.1.3}
H_i(X)\text{ es de grado $i$ en }X_0.
\end{equation}

De nuevo por (4.6.19),

\begin{equation}\tag{6.1.4}
\frac1p\,\psi_pR_S(i)\subset pR_S
\qquad\text{si }i\ne p.
\end{equation}

Como

\begin{equation}\tag{6.1.5}
H_i(X)\,R_S(j)\subset R_S(i+j),
\end{equation}

se tiene

\begin{equation}\tag{6.1.6}
\left(
\alpha_S-\frac1p\,\tau^{-1}\circ\psi_p\circ H_{p-1}(X)
\right)(R_S(1))
\subset pR_S.
\end{equation}

En consecuencia,

\begin{equation}\tag{6.1.7}
\left[
(\alpha_S)^{\operatorname{ord}(q)}
-\left(
\frac1p\,\tau^{-1}\circ\psi_p\circ H_{p-1}(X)
\right)^{\operatorname{ord}(q)}
\right]R_S(1)
\subset pR_S.
\end{equation}

Queda calcular $H_{p-1}(X)$. Por (4.4.4),

\pageoriginale

\begin{equation}\tag{6.1.8}
\begin{aligned}
H(X)
&=\prod_W\Theta(\pi A_WX_0X^W)\\
&=\prod_W\exp\!\left(
\sum_{r\geq0}\frac{(\pi A_WX_0X^W)^{p^r}}{p^r}
\right)\\
&=\prod_W\prod_{r\geq0}
\exp\!\left(\frac{(\pi A_WX_0X^W)^{p^r}}{p^r}\right)\\
&\equiv\prod_W\exp(\pi A_WX_0X^W)
\pmod{\text{términos de grado }\geq p}\\
&=\exp\!\left(\pi\sum_WA_WX_0X^W\right)
\pmod{\text{términos de grado }\geq p}\\
&=\exp(\pi X_0f(X))
\pmod{\text{términos de grado }\geq p}.
\end{aligned}
\end{equation}

Se sigue que

\begin{equation}\tag{6.1.9}
H_{p-1}(X)
=\frac{\pi^{p-1}X_0^{p-1}}{(p-1)!}\,f(X)^{p-1}.
\end{equation}

Se comprueba fácilmente que

\begin{equation}\tag{6.1.10}
\frac{\pi^{p-1}}p\equiv-1\pmod p,
\end{equation}

y, por tanto,

\begin{equation}\tag{6.1.11}
\frac1p\,H_{p-1}(X)
\equiv(X_0f(X))^{p-1}\pmod p.
\end{equation}

Esto termina la demostración de la Proposición~6.1.

\subsubsection*{Corolario 6.1.12.}
El operador lineal
\[
(\bar\alpha_S)^{\operatorname{ord}(q)}:W_S(1)\longrightarrow W_S(1)
\]
es precisamente
\[
\psi_q\circ(X_0f(X))^{q-1}.
\]

\pageoriginale

\noindent\emph{Demostración.}
Por (6.1.1),

\begin{equation*}
\begin{aligned}
(\bar\alpha_S)^{\operatorname{ord}(q)}
&=
\left(
\frac1p\,\tau^{-1}\circ\psi_p\circ(X_0f(X))^{p-1}
\right)^{\operatorname{ord}(q)}\\
&=
\psi_q\circ
\prod_{j=0}^{\sisi{\operatorname{ord}(a)}{\operatorname{ord}(q)}-1}
\left(X_0^{p^j}f^{\tau^j}(X^{p^j})\right)^{p-1}\\
&=
\psi_q\circ
\prod_{j=0}^{\operatorname{ord}(q)-1}
(X_0f(X))^{p^j(p-1)}\\
&=\psi_q\circ(X_0f(X))^{q-1}.
\end{aligned}
\end{equation*}

\subsubsection*{Corolario 6.1.13.}
Módulo $p$, se tiene la congruencia

\begin{equation*}
Z(t,V/\mathbb F_q)
\equiv
\frac{
\det\!\left(
1-t\,\psi_q\circ(X_0f(X))^{q-1}
\mathrel{\Big|}W_S(1)
\right)^{(-1)^{n+1}}
}{1-t}.
\end{equation*}

\subsection*{6.2.}
Queda calcular los grupos de cohomología de $\mathcal O_V$. Pongamos

\begin{equation}\tag{6.2.1}
\mathfrak F_q
=\text{la operación de elevar a la $q$-ésima potencia en $\mathcal O_V$}.
\end{equation}

Hay una sucesión exacta de haces sobre
$\mathbb P^{n+1}=\mathbb P$,

\begin{equation}\tag{6.2.2}
0\longrightarrow\mathcal O_{\mathbb P}(-d)
\xrightarrow{\,f\,}\mathcal O_{\mathbb P}
\longrightarrow\mathcal O_V\longrightarrow0,
\end{equation}

y el diagrama siguiente es conmutativo:

\begin{equation}\tag{6.2.3}
\begin{tikzcd}[column sep=large,row sep=large]
0 \arrow[r] &
\mathcal O_{\mathbb P}(-d) \arrow[r,"f"] \arrow[d,"f^{q-1}\mathfrak F_q"'] &
\mathcal O_{\mathbb P} \arrow[r] \arrow[d,"\mathfrak F_q"'] &
\mathcal O_V \arrow[r] \arrow[d,"\mathfrak F_q"'] &
0\\
0 \arrow[r] &
\mathcal O_{\mathbb P}(-d) \arrow[r,"f"] &
\mathcal O_{\mathbb P} \arrow[r] &
\mathcal O_V \arrow[r] &
0 .
\end{tikzcd}
\end{equation}

\pageoriginale

Teniendo en cuenta que

\begin{equation}\tag{6.2.4}
\begin{cases}
H^i(\mathbb P^{n+1},\mathcal O_{\mathbb P})=0,
&i\ne0,\\
H^i(\mathbb P^{n+1},\mathcal O_{\mathbb P}(-d))=0,
&i\ne0,n+1,
\end{cases}
\end{equation}

la sucesión exacta larga de cohomología asociada a (6.2.2) da

\begin{equation}\tag{6.2.5}
\begin{cases}
H^0(V,\mathcal O_V)\simeq\mathbb F_q,\\
H^i(V,\mathcal O_V)=0,&i\ne0,n,\\
H^n(V,\mathcal O_V)
\simeq H^{n+1}(\mathbb P^{n+1},\mathcal O_{\mathbb P}(-d)).
\end{cases}
\end{equation}

El diagrama conmutativo (6.2.3) da, por tanto,

\begin{equation}\tag{6.2.6}
\det\!\left(1-t\mathfrak F_q
\mathrel{\Big|}H^0(V,\mathcal O_V)\right)=1-t
\end{equation}

y

\begin{equation}\tag{6.2.7}
\det\!\left(1-t\mathfrak F_q
\mathrel{\Big|}H^n(V,\mathcal O_V)\right)
=
\det\!\left(
1-tf^{q-1}\mathfrak F_q
\mathrel{\Big|}
H^{n+1}(\mathbb P^{n+1},\mathcal O_{\mathbb P}(-d))
\right).
\end{equation}

Solo queda demostrar

\begin{equation}\tag{6.2.8}
\begin{aligned}
&\det\!\left(
1-tf^{q-1}\mathfrak F_q
\mathrel{\Big|}
H^{n+1}(\mathbb P^{n+1},\mathcal O_{\mathbb P}(-d))
\right)\\
&\qquad=
\det\!\left(
1-t\,\psi_q\circ(X_0f(X))^{q-1}
\mathrel{\Big|}W_S(1)
\right),
\end{aligned}
\end{equation}

lo cual es consecuencia de la proposición siguiente.

\subsubsection*{Proposición 6.3.}
Hay un isomorfismo

\begin{equation}\tag{6.3.1}
W_S(1)\simeq
\operatorname{Hom}_{\mathbb F_q}\!\left(
H^{n+1}(\mathbb P^{n+1},\mathcal O_{\mathbb P}(-d)),
\mathbb F_q\right)
\end{equation}

respecto al cual el operador

\begin{equation}\tag{6.3.2}
\psi_q\circ(X_0f(X))^{q-1}:W_S(1)\longrightarrow W_S(1)
\end{equation}

\pageoriginale

es el transpuesto de

\begin{equation}\tag{6.3.3}
f(X)^{q-1}\mathfrak F_q:
H^{n+1}(\mathbb P^{n+1},\mathcal O_{\mathbb P}(-d))
\longrightarrow
H^{n+1}(\mathbb P^{n+1},\mathcal O_{\mathbb P}(-d)).
\end{equation}

\noindent\emph{Demostración.}
El cálculo de
$H^{n+1}(\mathbb P^{n+1},\mathcal O_{\mathbb P}(-d))$ mediante el
recubrimiento estándar de $\mathbb P^{n+1}$ da

\begin{equation}\tag{6.3.4}
\begin{aligned}
H^{\sisi{i+1}{n+1}}(\mathbb P^{n+1},\mathcal O_{\mathbb P}(-d))
&=\text{el espacio vectorial sobre $\mathbb F_q$ con base los monomios}\\
&\quad X_1^{W_1}\cdots X_{n+2}^{W_{n+2}},
\qquad W_i\in\mathbb Z,\quad\sum_iW_i=-d,\\
&\sisi{=}{\mathrel{\big/}}\text{el subespacio sobre $\mathbb F_q$ con base aquellos monomios}\\
&\quad X_1^{W_1}\cdots X_{n+2}^{W_{n+2}}
\text{ para los que existe }i\text{ con }W_i\geq0.
\end{aligned}
\end{equation}

Por otra parte, $W_S(1)=R_S(1)\otimes_{\mathcal O}\mathbb F_q$, y

\begin{equation}\tag{6.3.5}
\begin{aligned}
R_S(1)
&=\text{el módulo libre sobre $\mathcal O$ con base los monomios}\\
&\quad
(\pi X_0)X_1^{U_1}\cdots X_{n+2}^{U_{n+2}},
\qquad
\sum_iU_i=d,\quad U_i>0\ \forall i.
\end{aligned}
\end{equation}

El emparejamiento

\begin{equation}\tag{6.3.6}
(\ ,\ ):
W_S(1)\times
H^{n+1}(\mathbb P^{n+1},\mathcal O_{\mathbb P}(-d))
\longrightarrow\mathbb F_q
\end{equation}

definido por

\begin{equation}\tag{6.3.7}
\left(
\operatorname{clase}\!\left(
\pi X_0X_1^{U_1}\cdots X_{n+2}^{U_{n+2}}
\right),
X_1^{W_1}\cdots X_{n+2}^{W_{n+2}}
\right)
=
\begin{cases}
1,&U_i=-W_i\text{ para todo }i,\\
0,&\text{en otro caso},
\end{cases}
\end{equation}

es claramente una dualidad perfecta de espacios vectoriales sobre
$\mathbb F_q$. Se comprueba de inmediato que
$\psi_q\circ(X_0f(X))^{q-1}$ es el transpuesto de
$f(X)^{q-1}\mathfrak F_q$, lo cual demuestra la Proposición~6.3.

\pageoriginale

\section*{Bibliografía}

\begin{enumerate}[label={\textup{[\arabic*]}},leftmargin=2.4em]
\item J.~Ax, «Zeroes of Polynomials Over Finite Fields»,
  \emph{Amer. J. Math.} \textbf{86} (1964), 255--261.
\item B.~Dwork, «On the Rationality of the Zeta Function of an Algebraic
  Variety», \emph{Amer. J. Math.} \textbf{82} (1960), 631--648.
\item B.~Dwork, «An Algebraic Variety», en \emph{Proceedings of the
  Woods Hole Summer Institute 1964} (mimeografiado).
\item B.~Dwork, «On the Zeta Function of a Hypersurface II»,
  \emph{Ann. of Math.} (2) \textbf{80} (1964), 227--299.
\item A.~Grothendieck, «Le Groupe de Brauer III», en
  \emph{Dix Exposés sur la Cohomologie des Schémas},
  North-Holland Publishing Co., 1968.
\item J.-P.~Serre, \emph{Groupes algébriques et corps de classes},
  Hermann, París, 1959.
\item J.-P.~Serre, «Endomorphismes complètement continus des espaces de
  Banach $p$-adiques», \emph{Publ. Math. I.H.E.S.} n.º~12 (1962).
\item C.~Chevalley, «Démonstration d'une hypothèse de M.~Artin»,
  \emph{Abh. Math. Sem. Univ. Hamburg} \textbf{11} (1936), 73--75.
\end{enumerate}
